% AA-47 | Algorytmy Rischa i Kovacica — zakres zatwierdzony przez kontynuację.
% Dowody pełne: integracja C(x), redukcja Hermite'a, test Rothsteina--Tragera,
% r'+hr=s przy h wielomianowym, wymierna galaz Riccatiego dla q wielomianowego.
% Wyniki cytowane: pelny algorytm Rischa i pelny algorytm Kovacica (przypadki 2, 3).
\chapter{Algorytmy Rischa i Kovacica}
\chaptermark{Algorytmy Rischa i Kovacica}
\label{aa:ch:risch-kovacic}

Twierdzenie Liouville'a z rozdziału
\ref{aa:ch:liouville-elementary-integration}
daje kryterium istnienia elementarnej funkcji pierwotnej,
a teoria Picarda--Vessiota z rozdziału
\ref{aa:ch:picard-vessiot}
wiąże rozwiązania liouvillowskie równań liniowych
z grupami algebraicznymi. Są to twierdzenia
strukturalne. W tym rozdziale przejdziemy
od istnienia do \emph{rozstrzygalnych obliczeń}.

Dla przejrzystości rozdzielamy dwa poziomy:
\emph{pełne dowody} procedur nad
$\mathbb C(x)$, dla równań $r'+hr=s$
z $h\in\mathbb C[x]$, $s\in\mathbb C(x)$,
oraz dla wymiernego równania Riccatiego
z wielomianowym współczynnikiem;
\emph{klasyczne wyniki cytowane}
dotyczą ogólnego algorytmu Rischa
i pozostałych gałęzi algorytmu Kovacica.
Nie utożsamiamy pierwszego zakresu
z pełną implementacją drugiego.

\section{Model obliczeń i znaczenie odpowiedzi algorytmu}

\begin{definicja}[Efektywnie zadane współczynniki]
\label{aa:def:effective-constants}
Przy obliczeniach nad ciałem $C$ charakterystyki
$0$ zakładamy dostęp do dokładnych działań,
testu równości i efektywnego rozkładu
wielomianów używanych w procedurze
na czynniki w odpowiednim skończonym
algebraicznym rozszerzeniu ciała współczynników.
Jako podstawowy przykład można przyjąć
ciało liczb algebraicznych
$\overline{\mathbb Q}\subset\mathbb C$.
Twierdzenia matematyczne poniżej formułujemy
nad $\mathbb C$, a słowo \emph{algorytm}
oznacza skończoną liczbę wskazanych działań
na efektywnie reprezentowanych danych.
\end{definicja}

Przy zadaniu współczynników przez dowolne
nienazwane liczby zespolone bez testu równości
nie otrzymujemy modelu efektywnego.
W szczególności \emph{pełne}
twierdzenia o rozstrzygalności wymagają
jawnego określenia sposobu reprezentacji
ciała współczynników.

\section{Całkowanie wymierne: pełny algorytm}

\begin{twierdzenie}[Rozkład na część wymierną i logarytmy]
\label{aa:tw:rational-integration-algorithm}
Niech $f\in\mathbb C(x)$. Istnieją
$R\in\mathbb C(x)$, parami różne
$a_1,\ldots,a_m\in\mathbb C$
i $c_1,\ldots,c_m\in\mathbb C$ takie, że
\begin{equation}\label{aa:eq:rational-integration}
 f=R'+\sum_{j=1}^{m}\frac{c_j}{x-a_j}.
\end{equation}
Można je wyznaczyć w skończonej liczbie
działań arytmetycznych i rozkładów wielomianów.
Ponadto $f$ ma funkcję pierwotną
w $\mathbb C(x)$ wtedy i tylko wtedy,
gdy wszystkie residua $c_j$ są zerowe.
Niezależnie od tego $f$ ma lokalną
elementarną funkcję pierwotną.
\end{twierdzenie}
\begin{proof}
Dzielimy licznik przez mianownik,
a następnie rozkładamy część właściwą
na ułamki proste:
\[
 f=P(x)+
 \sum_{j=1}^{m}\sum_{k=1}^{n_j}
         \frac{b_{j,k}}{(x-a_j)^k}.
\]
Jest to skończona procedura nad
ciałem, w którym rozkłada się mianownik.
Wielomian $P$ całkujemy wyraz po wyrazie.
Dla $k\ge2$ element
$-b_{j,k}(x-a_j)^{1-k}/(k-1)$
ma pochodną
$b_{j,k}(x-a_j)^{-k}$.
Po odjęciu otrzymanych pochodnych
pozostają tylko składniki
$b_{j,1}/(x-a_j)$.
Daje to~\eqref{aa:eq:rational-integration}
z $c_j=b_{j,1}$.
W każdym skończonym punkcie
residuum pochodnej funkcji wymiernej
jest zerowe: pochodna szeregu
Laurenta $\sum_k v_k(x-a)^k$
nie ma współczynnika przy
$(x-a)^{-1}$, gdyż dla $k=0$
otrzymujemy $k v_k=0$.
Jeśli więc $f$ ma wymierną funkcję
pierwotną, wszystkie $c_j$ muszą być
zerowe. Odwrotnie, przy znikających
residuach $f=R'$, więc taka
funkcja pierwotna istnieje.
W ogólnym przypadku na ustalonych
lokalnych gałęziach można przyjąć
$R+\sum_jc_j\log(x-a_j)$.
\end{proof}

\begin{definicja}[Redukcja Hermite'a]
\label{aa:def:hermite-rational}
\emph{Redukcją Hermite'a funkcji wymiernej}
nazywamy konstrukcję, która odejmuje
pochodne wymiernych części głównych
rzędu co najmniej dwa i sprowadza
$f$ do sumy pochodnej wymiernej
oraz funkcji wymiernej mającej
wyłącznie bieguny proste.
W dowodzie
Twierdzenia~\ref{aa:tw:rational-integration-algorithm}
wykonaliśmy taką redukcję
przez rozkład na ułamki proste.
Warianty bez rozdzielania mianownika
na czynniki liniowe używają
dzielenia wielomianów i ich NWD.
\end{definicja}

\begin{przyklad}[Obliczenie obu części]
\label{aa:prz:hermite-example}
Dla
$f=2x+3/(x-1)^3+5/(x-1)$
otrzymujemy
\[
 R=x^2-\frac{3}{2(x-1)^2},
 \qquad
 f=R'+\frac5{x-1}.
\]
Zatem na dowolnym obszarze,
na którym można wybrać gałąź
logarytmu $x-1$, funkcją
pierwotną jest
$R+5\log(x-1)$.
Nie istnieje wymierna funkcja
pierwotna, ponieważ residuum
w $1$ wynosi $5$.
\end{przyklad}

\begin{twierdzenie}[Wykorzystanie wypadkowej do residuów]
\label{aa:tw:rothstein-trager}
Niech $A,B\in\mathbb C[x]$,
$\deg A<\deg B$, a $B$ będzie
unormowanym wielomianem bez
pierwiastków wielokrotnych.
Dla każdego pierwiastka $a$ wielomianu
$B$ residuum $A/B$ jest równe
$c_a=A(a)/B'(a)$.
Wielomian
\begin{equation}\label{aa:eq:rt-resultant}
 T(z)=\operatorname{Res}_x(B(x),A(x)-zB'(x))
\end{equation}
ma (z uwzględnieniem krotności)
dokładnie te wartości $c_a$ jako pierwiastki.
Jeśli $c$ przebiega różne pierwiastki $T$,
a $P_c=\gcd(B,A-cB')$ jest unormowany,
to zachodzi tożsamość
\begin{equation}\label{aa:eq:rt-reconstruction}
 \frac{A}{B}=\sum_{c:T(c)=0}
              c\frac{P_c'}{P_c}.
\end{equation}
Zatem przy efektywnym rozkładzie
$T$ i obliczaniu NWD można znaleźć
część logarytmiczną bez wcześniejszego
rozkładu $B$ na czynniki liniowe.
\end{twierdzenie}
\begin{proof}
Ponieważ pierwiastki $a$ wielomianu $B$
są proste, lokalnie
$B(x)=(x-a)B'(a)+O((x-a)^2)$,
co daje podane residuum.
Własność wypadkowej dla
unormowanego $B$ daje
\[
 T(z)=\prod_{B(a)=0}
           (A(a)-zB'(a))
     =\left(\prod_{B(a)=0}B'(a)\right)
       \prod_{B(a)=0}(c_a-z).
\]
Zatem pierwiastki wypadkowej są
dokładnie wartościami residuów.
Pierwiastkami $P_c$ są właśnie
te pierwiastki $a$ wielomianu $B$,
dla których $c_a=c$. W szczególności
\[
 \frac{P_c'}{P_c}
 =\sum_{\substack{B(a)=0\\c_a=c}}
             \frac1{x-a}.
\]
Obie strony~\eqref{aa:eq:rt-reconstruction}
mają te same części główne
przy każdym skończonym biegunie;
obie są właściwymi funkcjami wymiernymi
(przy nieskończoności mają granicę
zerową). Ich różnica jest wielomianem
bez części wielomianowej, a więc zerem.
\end{proof}

\begin{przyklad}[Wypadkowa bez rozkładu mianownika]
\label{aa:prz:rt-example}
Dla $A=1$, $B=x^2+1$ mamy
$B'=2x$ i
\[
 T(z)=\operatorname{Res}_x(x^2+1,1-2zx)
     =1+4z^2.
\]
Pierwiastkami są $c=\pm i/2$.
Odpowiednie NWD wybierają czynniki
$x+i$ i $x-i$. Otrzymujemy
\[
 \frac1{x^2+1}
 =\frac1{2i}\left(\frac1{x-i}
                         -\frac1{x+i}\right).
\]
Dla $A=2x$ otrzymalibyśmy
jedno residuum $c=1$, a
$P_1=x^2+1$, skąd
$2x/(x^2+1)=B'/B$.
\end{przyklad}

\section{Efektywna część algorytmu Rischa: krok wykładniczy}

Niech $h\in\mathbb C[x]$ i $s\in\mathbb C(x)$.
Dla niezerowego $h$ wielomian
$p(x)=\int h(x)\,\mathrm dx$
jest niestały; na mocy
Lematu~\ref{aa:lem:exp-polynomial-constant-field}
element $t=e^{p(x)}$ jest
przestępny nad $\mathbb C(x)$
i nie dodaje stałych.
Z Twierdzenia
\ref{aa:tw:hyperexponential-criterion}
wynika, że element
$st$ ma elementarną funkcję pierwotną
wtedy i tylko wtedy, gdy
\begin{equation}\label{aa:eq:risch-first-order}
 r'+hr=s
\end{equation}
ma rozwiązanie $r\in\mathbb C(x)$.
Pokażemy, że w tym zakresie
test jest faktycznie \emph{skończonym
algorytmem decyzyjnym}.

\begin{lemat}[Bieguny wymiernego rozwiązania]
\label{aa:lem:risch-poles}
Niech $h\in\mathbb C[x]$ i
$s,r\in\mathbb C(x)$ spełniają
$r'+hr=s$. Jeżeli $r$ ma
w skończonym punkcie $a$
biegun rzędu $j\ge1$, to
$s$ ma w $a$ biegun dokładnie
rzędu $j+1$. Współczynnik
wiodący części głównej $s$
jest równy $-j$ razy
współczynnik wiodący
części głównej $r$.
W szczególności rozwiązanie
$r$ nie może mieć biegunów
w punktach regularności $s$,
a prosty biegun $s$ uniemożliwia
istnienie rozwiązania.
\end{lemat}
\begin{proof}
W lokalnej zmiennej $z=x-a$
zapiszmy $r=b_{-j}z^{-j}+
O(z^{1-j})$, gdzie $b_{-j}\ne0$.
Wtedy $r'=-j b_{-j}z^{-j-1}
+O(z^{-j})$, natomiast
$hr$ ma biegun rzędu
co najwyżej $j$, ponieważ
$h$ jest regularny w $a$.
Składnik o rzędzie $j+1$
nie może się skasować.
Zatem $s$ ma dokładnie
taki biegun i podany
współczynnik.
\end{proof}

\begin{twierdzenie}[Algorytm wymiernego równania Rischa]
\label{aa:tw:risch-polynomial-h-algorithm}
Niech $h\in\mathbb C[x]\setminus\{0\}$,
$s\in\mathbb C(x)$.
Istnieje skończona procedura,
oparta na rozkładzie mianownika
$s$ oraz rozwiązywaniu
skończonych układów równań liniowych,
która rozstrzyga istnienie
$r\in\mathbb C(x)$ spełniającego
$r'+hr=s$ i, jeśli rozwiązanie
istnieje, wyznacza je.
Dla $h=0$ kryterium
istnienia wymiernej funkcji pierwotnej
jest podane w
Twierdzeniu~\ref{aa:tw:rational-integration-algorithm}.
\end{twierdzenie}
\begin{proof}
Niech $a$ będzie skończonym biegunem
$s$ rzędu $m$. Jeżeli $m=1$,
Lemat~\ref{aa:lem:risch-poles}
wyklucza rozwiązanie.
Gdy $m\ge2$, ewentualne rozwiązanie
ma w $a$ część główną
$R_a(z)=\sum_{j=1}^{m-1}b_jz^{-j}$.
Podstawiamy ją do $r'+hr=s$
i porównujemy kolejno współczynniki
przy $z^{-m},z^{-m+1},\ldots,z^{-2}$.
Układ jest trójkątny:
współczynnik przy $z^{-j-1}$
zawiera $-j b_j$ oraz wyrazy
zależące wyłącznie od
$b_{j+1},\ldots,b_{m-1}$,
bo $h$ ma nieujemne potęgi $z$.
Ponieważ charakterystyka wynosi
$0$, wszystkie $b_j$ są
jednoznacznie wyznaczone.
Pozostaje sprawdzić współczynnik
przy $z^{-1}$. Jeżeli
różni się od odpowiedniego
współczynnika $s$, nie ma
rozwiązania. Jeśli się zgadza,
wyznaczona część główna usuwa
wszystkie bieguny w $a$.
Powtarzamy to dla każdego
skończonego bieguna $s$.
Otrzymujemy wymierną funkcję
$R=\sum_aR_a(x-a)$ i wielomian
$S=s-R'-hR\in\mathbb C[x]$.

Gdyby pierwotne równanie miało
rozwiązanie wymierne $r$,
to Lemat~\ref{aa:lem:risch-poles}
i jednoznaczność porównań
współczynników wymuszają,
że $r-R$ nie ma żadnych
skończonych biegunów.
Zatem $r-R=P\in\mathbb C[x]$.
Pozostaje zdecydować, czy
$P'+hP=S$ ma rozwiązanie wielomianowe.

Jeżeli $h$ ma dodatni stopień
$d$ i $S\ne0$, to
dla niezerowego $P$ mamy
$\deg(P'+hP)=\deg P+d$.
Wobec tego stopień $P$
musi być równy
$\deg S-d$. Jeśli liczba
ta jest ujemna, rozwiązania
nie ma; w przeciwnym razie
podstawiamy wielomian
$P=\sum_{k=0}^{\deg S-d}p_kx^k$
i rozstrzygamy skończony
\emph{liniowy} układ równań
na współczynniki $p_k$.
Jeśli $S=0$, rozwiązaniem
jest $P=0$, a innego
wielomianowego rozwiązania
jednorodnego nie ma.

Jeśli $h=c\in\mathbb C^\times$
jest stały, operator
$P\mapsto P'+cP$ jest
odwracalny na każdym
skończenie wymiarowym
zbiorze wielomianów
stopnia co najwyżej $\deg S$.
Można to zobaczyć bezpośrednio:
\[
 P=\sum_{k=0}^{\deg S}
                 \frac{(-1)^k S^{(k)}}{c^{k+1}}
\]
spełnia $P'+cP=S$, ponieważ
suma teleskopuje, a wyższe
pochodne $S$ znikają.
Dla $S=0$ przyjmujemy $P=0$.
We wszystkich przypadkach
wykonaliśmy skończoną liczbę
działań. Uzasadniliśmy też
konieczność i dostateczność
każdego testu.
\end{proof}

\begin{przyklad}[Obliczenie zakończone sukcesem]
\label{aa:prz:risch-e-success}
Dla $h=1$ i $s=1/x^2-1/x$
jedyny możliwy biegun $r$
jest prosty i znajduje się
w $0$. Z wyrazu $x^{-2}$
wynika $R=-1/x$.
Sprawdzamy
$R'+R=1/x^2-1/x=s$,
zatem $S=0$.
Wynik algorytmu jest
\[
 \int e^x\left(\frac1{x^2}-\frac1x\right)
 \,\mathrm dx=-\frac{e^x}{x}+C.
\]
Dla $h=1$, $s=1/x$
procedura kończy się negatywnie
od razu po wykryciu bieguna prostego.
\end{przyklad}

\begin{przyklad}[Całka Gaussa jako odpowiedź negatywna]
\label{aa:prz:risch-gauss-decision}
Dla $h=-2x$ i $s=1$
nie ma biegunów skończonych.
Pozostało $P'-2xP=1$.
Ponieważ $\deg S=0<1=\deg h$,
algorytm dowodzi nieistnienia
wymiernego $P$.
W połączeniu z kryterium
\ref{aa:tw:hyperexponential-criterion}
jest to dowód braku
elementarnej funkcji pierwotnej
$e^{-x^2}$, nie tylko niepowodzenie
heurystycznego programu.
\end{przyklad}

\begin{uwaga}[Pełny algorytm Rischa -- wynik klasyczny]
\label{aa:uw:full-risch-scope}
Właściwy algorytm Rischa
wychodzi daleko poza pole
$\mathbb C(x,e^{p(x)})$:
przetwarza odpowiednio
reprezentowane wieże
rozszerzeń elementarnych,
redukuje części algebraiczne
i przestępne i rozwiązuje
równania różniczkowe o bardziej
złożonych współczynnikach.
Jego zupełność w pełnym
zakresie przyjmujemy jako
wynik klasyczny (Risch;
Bronstein, \emph{Symbolic Integration I}).
W tym rozdziale udowodniono
kompletną procedurę \emph{tylko}
dla opisanej klasy wymiernej
i dla $h\in\mathbb C[x]$.
Nie twierdzimy, że
implementacje systemów CAS
realizują wszystkie przypadki
ogólnego twierdzenia.
\end{uwaga}

\section{Algorytm Kovacica: redukcja i klasyfikacja}

\begin{lemat}[Usunięcie pierwszej pochodnej]
\label{aa:lem:kovacic-normal-form}
Niech $a,b\in\mathbb C(x)$.
Lokalnie, w punkcie bez osobliwości
odpowiednich funkcji, podstawienie
$y=z\exp(-\tfrac12\int a(x)\,\mathrm dx)$
przekształca równanie
\[
 y''+a(x)y'+b(x)y=0
\]
w równanie $z''=q(x)z$, gdzie
\begin{equation}\label{aa:eq:kovacic-normal-q}
 q=\frac{a'}2+\frac{a^2}4-b
 \in\mathbb C(x).
\end{equation}
Podstawienie jest odwracalne
w lokalnym rozszerzeniu
liouvillowskim.
\end{lemat}
\begin{proof}
Niech $g=\exp(-\tfrac12\int a\,\mathrm dx)$;
wtedy $g'/g=-a/2$.
Dla $y=gz$ mamy
$y'=g(z'-az/2)$ oraz
$y''=g(z''-az'+(a^2/4-a'/2)z)$.
Po podstawieniu w równanie
znikają wyrazy przy $z'$,
a pozostaje
$z''+(b-a'/2-a^2/4)z=0$,
co jest równoważne podanej tezie.
Współczynnik $g$ jest niezerowy,
zatem podstawienie jest odwracalne.
\end{proof}

\begin{lemat}[Równanie Riccatiego]
\label{aa:lem:kovacic-riccati}
Niech $q\in\mathbb C(x)$ i
$z$ będzie niezerowym
lokalnym rozwiązaniem
$z''=qz$. Wówczas
$r=z'/z$ spełnia
$r'+r^2=q$.
Odwrotnie, jeśli element
$r$ dowolnego rozszerzenia
różniczkowego spełnia
$r'+r^2=q$, a
$z\ne0$ spełnia $z'/z=r$,
to $z''=qz$.
\end{lemat}
\begin{proof}
Z reguły ilorazu
$r'=z''/z-(z'/z)^2
=q-r^2$.
Odwrotnie, różniczkując
$z'=rz$ otrzymujemy
$z''=(r'+r^2)z=qz$.
\end{proof}

\begin{twierdzenie}[Trzy przypadki Kovacica -- wynik klasyczny]
\label{aa:tw:kovacic-classification}
Niech $q\in\mathbb C(x)$,
a współczynniki będą zadane
efektywnie. Dla równania
$z''=qz$ klasyczny algorytm
Kovacica rozstrzyga, czy istnieje
niezerowe rozwiązanie
liouvillowskie, i jeśli istnieje,
konstruuje takie rozwiązanie.
Wynik klasyfikacyjny dla
logarytmicznej pochodnej
$r=z'/z$ dopuszcza następujące
rodzaje przypadków:
\begin{enumerate}
\item $r\in\mathbb C(x)$
(przypadek redukowalny);
\item $r$ jest algebraiczne
nad $\mathbb C(x)$ stopnia
$2$ (przypadek grupy imprymitywnej
grupy Galois);
\item $r$ jest algebraiczne
stopnia $4$, $6$ lub $12$
(przypadek skończonej
prymitywnej grupy Galois);
\item nie ma niezerowego
rozwiązania liouvillowskiego.
\end{enumerate}
W przypadkach 1--3
algorytm wyszukuje skończoną
liczbę możliwości wynikających
z rozwinięć lokalnych
i sprawdza tożsamości
wymierne lub algebraiczne.
\end{twierdzenie}

\begin{uwaga}[Źródło i granice klasyfikacji]
\label{aa:uw:kovacic-citation}
Powyższy wynik cytujemy z:
J.~Kovacic, \emph{An algorithm
for solving second order linear
homogeneous differential equations},
\emph{Journal of Symbolic Computation}
2 (1986), 3--43.
Klasyfikacja korzysta z
podgrup algebraicznych
$\mathrm{SL}_2(\mathbb C)$
oraz z teorii Picarda--Vessiota
w rozdziale~\ref{aa:ch:picard-vessiot}.
Nie przedstawiamy rozgałęzień
2 i 3 jako samodzielnie
udowodnionych w tym rozdziale,
ani nie traktujemy testu
gałęzi 1 jako testu
liouvillowskości w całej klasie.
\end{uwaga}

\section{Gałąź wymierna Riccatiego dla wielomianowego współczynnika}

\begin{lemat}[Bieguny wymiernego rozwiązania Riccatiego]
\label{aa:lem:kovacic-rational-poles}
Niech $q\in\mathbb C[x]$
i niech $r\in\mathbb C(x)$
spełnia $r'+r^2=q$.
Każdy skończony biegun $r$
jest prosty i ma residuum $1$.
Istnieją więc
$\omega\in\mathbb C[x]$
oraz unormowany, niezerowy
wielomian $P\in\mathbb C[x]$
bez pierwiastków wielokrotnych
takie, że
\begin{equation}\label{aa:eq:kovacic-rational-r}
 r=\omega+\frac{P'}P.
\end{equation}
Co więcej, $P$ ma dokładnie
wszystkie skończone bieguny $r$
jako pierwiastki.
\end{lemat}
\begin{proof}
Jeśli $r$ ma biegun rzędu
$m\ge2$ w $a$, to $r^2$
ma biegun rzędu $2m$,
a $r'$ rzędu $m+1<2m$,
co przeczy regularności $q$.
Dla bieguna prostego piszemy
$r=c(x-a)^{-1}+O(1)$,
$c\ne0$. Współczynnik
przy $(x-a)^{-2}$ w
$r'+r^2$ wynosi
$c^2-c$; musi być zerowy,
więc $c=1$.
Odejmując od $r$
sumę $1/(x-a)$ po jego
skończonych biegunach
otrzymujemy funkcję wymierną
bez biegunów skończonych,
a więc wielomian $\omega$.
Jeśli $P$ jest iloczynem
$(x-a)$ po tych biegunach
(jeśli ich nie ma, $P=1$),
to $P'/P=\sum_a1/(x-a)$,
co daje tezę.
\end{proof}

\begin{twierdzenie}[Pełny test wymiernego Riccatiego dla $q$ wielomianowego]
\label{aa:tw:kovacic-polynomial-rational-decision}
Niech $q\in\mathbb C[x]$
będzie zadany efektywnie.
Istnieje skończony algorytm,
który rozstrzyga, czy
równanie
\begin{equation}\label{aa:eq:kovacic-poly-riccati}
 r'+r^2=q
\end{equation}
ma rozwiązanie $r\in\mathbb C(x)$.
Jeżeli $q$ jest stałe,
rozwiązaniem jest dowolny
stały pierwiastek kwadratowy $q$.
Jeżeli $\deg q$ jest nieparzysty,
nie ma rozwiązania wymiernego.
Gdy $\deg q=2d\ge2$,
procedura obejmuje dokładnie
dwa możliwe wielomiany
$\omega$ stopnia $d$,
obliczenie jednego wymuszonego
stopnia $n$ dla każdego
z nich, oraz rozwiązanie
skończonego układu
\emph{liniowego} na
współczynniki wielomianu $P$
stopnia $n$.
\end{twierdzenie}
\begin{proof}
W Lemacie~\ref{aa:lem:kovacic-rational-poles}
otrzymaliśmy postać
$r=\omega+P'/P$, gdzie
$\omega\in\mathbb C[x]$
i $P$ jest unormowany.
Podstawienie do Riccatiego daje
\begin{equation}\label{aa:eq:kovacic-polynomial-P}
 P''+2\omega P'
 +(\omega'+\omega^2-q)P=0.
\end{equation}
Jeżeli $\deg P=n$, to
$P'/P=n/x+O(x^{-2})$
w nieskończoności.
Jeśli $\omega$ ma dodatni
stopień $d$, to człony
$\omega'$ oraz
$2\omega(P'/P)$
mają stopień co najwyżej
$d-1$ w części wielomianowej,
a pochodna $P'/P$ i jego
kwadrat mają wyłącznie
ujemne potęgi.
Stąd wielomianowa część
$r'+r^2$ ma stopień $2d$
i współczynnik wiodący
równy kwadratowi
współczynnika wiodącego
$\omega$. W szczególności
$q$ musi mieć dodatni
\emph{parzysty} stopień.
Jeśli $\omega$ jest stały,
wielomianowa część
$r'+r^2$ jest stała,
więc nie otrzymamy
$q$ dodatniego stopnia.
To dowodzi przeszkody
dla nieparzystego stopnia.

Załóżmy teraz $\deg q=2d\ge2$.
Niech $a_d$ oznacza
współczynnik wiodący
$\omega$. Musi zachodzić
$a_d^2=[x^{2d}]q$, są
zatem dwa możliwe wybory
$a_d\in\mathbb C^\times$.
Dla każdego wybieramy
kolejno współczynniki
$a_{d-1},\ldots,a_0$
tak, aby
\[
 \deg(q-\omega^2)\le d-1.
\]
Współczynniki przy
$x^{2d-1},\ldots,x^d$
wyznaczają rekurencyjnie
wszystkie $a_j$,
ponieważ współczynnik
przy nieznanym $a_j$
wynosi $2a_d\ne0$.
Są to dokładnie
dwa możliwe wielomiany
$\omega$.

Porównajmy teraz
współczynnik przy
$x^{d-1}$ w tożsamości
\[
 q-\omega^2-\omega'
 =2\omega\frac{P'}P+
 \left(\frac{P'}P\right)'
 +\left(\frac{P'}P\right)^2.
\]
Ponieważ $P'/P=n/x+O(x^{-2})$,
jego prawa strona ma
współczynnik przy $x^{d-1}$
równy $2a_dn$. Dlatego
\begin{equation}\label{aa:eq:kovacic-degree-n}
 n=\frac{[x^{d-1}](q-\omega^2-\omega')}
          {2a_d}.
\end{equation}
Jeśli otrzymana liczba
nie należy do $\mathbb Z_{\ge0}$,
to wybór $\omega$ jest
niemożliwy.

W pozostałych przypadkach
podstawiamy
$P=x^n+\sum_{j=0}^{n-1}p_jx^j$
do~\eqref{aa:eq:kovacic-polynomial-P}.
Porównanie współczynników
daje skończony
\emph{liniowy} układ równań
dla $p_0,\ldots,p_{n-1}$.
Każde jego rozwiązanie
wyznacza wielomian $P$
spełniający tożsamość.
Nie może on mieć
wielokrotnych pierwiastków:
gdyby $a$ był pierwiastkiem
krotności $m\ge2$, wyraz
$P''$ w~\eqref{aa:eq:kovacic-polynomial-P}
miałby w $a$ rząd
$m-2$ z niezerowym współczynnikiem,
podczas gdy pozostałe
wyrazy miałyby rząd
co najmniej $m-1$.
Stąd konstrukcja $r$
z~\eqref{aa:eq:kovacic-rational-r}
jest poprawna i spełnia
\eqref{aa:eq:kovacic-poly-riccati}.

Na odwrót każde wymierne
rozwiązanie miało postać
z Lematu~\ref{aa:lem:kovacic-rational-poles},
a porównanie współczynników
wymuszało jeden z dwóch
wyborów $\omega$, ten sam
stopień $n$ i rozwiązanie
jednego z badanych układów.
Algorytm nie odrzuci więc
żadnego wymiernego rozwiązania.
Dla $q$ stałego zawsze
wystarcza stałe $r$
z $r^2=q$, co kończy dowód.
\end{proof}

\begin{przyklad}[Dwa wyniki pozytywne]
\label{aa:prz:kovacic-polynomial-success}
Dla $q=x^2+1$ możemy
wybrać $\omega=x$ i
\eqref{aa:eq:kovacic-degree-n}
daje $n=0$; wielomian
$P=1$ spełnia
\eqref{aa:eq:kovacic-polynomial-P}.
Otrzymujemy $r=x$ oraz
rozwiązanie
$z=e^{x^2/2}$.
Dla $q=x^2+3$ ten sam
wybór $\omega=x$
daje $n=1$, a
$P=x$ spełnia
$P''+2xP'+(1+x^2-(x^2+3))P=0$.
Stąd
\[
 r=x+\frac1x,\qquad
 z=xe^{x^2/2},\qquad
 z''=(x^2+3)z.
\]
Obydwa wyniki można
zweryfikować niezależnie
przez różniczkowanie.
\end{przyklad}

\begin{przyklad}[Dlaczego gałąź wymierna nie jest całą procedurą]
\label{aa:prz:kovacic-quadratic-example}
Rozważmy
$q=x+5/(16x^2)$.
W nieskończoności $q$
ma część wielomianową
stopnia nieparzystego,
więc równanie Riccatiego
nie ma \emph{wymiernego}
rozwiązania: gdyby miało,
wielomianowa część
$r'+r^2$ byłaby stopnia
parzystego albo stała.
Jednakże
\[
 r=\sqrt{x}-\frac1{4x}
\]
spełnia $r'+r^2=q$,
co sprawdzamy bezpośrednio:
składniki z $x^{-1/2}$
znoszą się, a przy
$x^{-2}$ zostaje $5/16$.
Jest to element
stopnia $2$ nad
$\mathbb C(x)$. Odpowiada
mu rozwiązanie
\[
 z=x^{-1/4}
       \exp\left(\frac23 x^{3/2}\right)
\]
na lokalnie wybranych
gałęziach. Dlatego
odrzucenie pierwszej gałęzi
nie jest dowodem braku
rozwiązań liouvillowskich.
\end{przyklad}

\begin{przyklad}[Równanie Airy'ego]
\label{aa:prz:kovacic-airy}
Dla $q=x$ mamy
$\deg q=1$, zatem test
wymierny daje odpowiedź negatywną.
Silniejsze stwierdzenie,
że równanie Airy'ego
$z''=xz$ \emph{nie ma żadnego}
niezerowego rozwiązania
liouvillowskiego,
udowodniliśmy już w
Twierdzeniu~\ref{aa:tw:pv-airy-group},
korzystając z kryterium
Picarda--Vessiota i
wykluczenia wszystkich
\emph{algebraicznych}
rozwiązań Riccatiego.
Nie przypisujemy tej
silniejszej konkluzji
samemu testowi wymiernemu.
\end{przyklad}

\section{Zestawienie i granice rozstrzygalności}

\begin{itemize}
\item Dla $f\in\mathbb C(x)$
procedura wymierna
z~\ref{aa:tw:rational-integration-algorithm}
wyznacza pochodną wymierną,
residua i logarytmy, a
test znikania residuów
rozstrzyga istnienie
\emph{wymiernej} funkcji pierwotnej.
\item Dla $s\in\mathbb C(x)$,
$h\in\mathbb C[x]\setminus\{0\}$
procedura~\ref{aa:tw:risch-polynomial-h-algorithm}
i kryterium AA-45
rozstrzygają elementarność
funkcji pierwotnej
$s\exp(\int h\,\mathrm dx)$.
\item Dla $z''=qz$ z
$q\in\mathbb C[x]$
procedura~\ref{aa:tw:kovacic-polynomial-rational-decision}
rozstrzyga istnienie
\emph{wymiernego} rozwiązania
Riccatiego i daje wtedy
rozwiązanie liouvillowskie.
Ogólna procedura Kovacica
dla $q\in\mathbb C(x)$
obejmuje też rozwiązania
algebraiczne stopni
$2,4,6,12$; jej pełna
poprawność jest cytowana,
nie udowodniona.
\end{itemize}

W żadnym z tych problemów
numeryczne przybliżenie
funkcji pierwotnej albo
rozwiązania równania
różniczkowego nie zastępuje
dokładnej odpowiedzi o
elementarności lub
liouvillowskości.
